% Stage 8 restored fragment — 2.6.1 Compact circular aperture and singular-value decay (archive lines 7474-7642)
\subsection{Compact circular aperture and singular-value decay}
\label{sec:ch26.signature-aperture}
Before any eigenvalue formula is written, it helps to watch the phenomenon on a range. On a dish
compact range, place a feed above a circular scan disk on plane $z=0$ and sample tangential field on
$S_r$; on a spherical near-field scanner, place probes on a hull at $R\approx 2\lambda_0$ and
transform to far-zone coefficients \cite{hansen1988}. In either campaign, count how
many independent beam directions remain above the receiver noise floor. The answer depends on disk
or scan radius $a_{\mathrm{ap}}$ and source electrical size, not on how many spherical harmonics a
post-processor can index \cite{chu1948}. Electrical size and scan extent fix a
recoverable rank; singular values $\sigma_n$ of $\mathcal{O}_{\mathrm{meas}}\circ\mathcal{R}_\omega$
measure the radiation efficiency of each independent source organization through the aperture and
receiver chain, and fall off once the modal index exceeds what the geometry can resolve
\cite{stratton1941}. Synthesis harmonic caps therefore overstate physical rank when $\sigma_n$ has already fallen below the
noise floor---harmonics in a post-processor are labels, not independent export channels. Compact
decay of $\sigma_n(\mathcal{K})$ and the finite recoverable rank it defines are stated here once;
later subsections invoke Subsection~\ref{sec:ch26.signature-aperture} by reference rather than
repeating the measurement narrative \cite{harrington2001}.

Consider a circular scan aperture of radius \(a_{\mathrm{ap}}\) on plane \(z=0\) sampling the
far-zone tangential field from a compact source with \(k_0 a=1.5\). The measurement composition
$\mathcal{K}=\mathcal{O}_{\mathrm{meas}}\circ\mathcal{R}_\omega$ is compact on the source and
field spaces declared in Section~\ref{sec:ch213}; its singular values $\sigma_n$ rank radiation
efficiency of source organizations and decay, defining a practical recoverable rank and finite
effective dimension before any harmonic label is attached
\cite{stratton1941}.
\begin{figure}[t]
\centering
\ChOneFigGeo{%
\begin{tikzpicture}[ch1geo, scale=0.90]
 \draw[ch1shell] (0,0) circle (0.75);
 \node[ch1oncanvas, below=8mm] at (0,-0.85) {aperture radius $a_{\mathrm{ap}}$};
 \draw[ch1shellB, dashed] (0,0) circle (0.97);
 \node[ch1oncanvas, right=10mm] at (1.35,0) {$S_r$ far zone};
 \draw[->, thick, gray] (0.6,0.2) -- (1.2,0.35);
 \node[ch1oncanvas, font=\scriptsize] at (0.9,0.55) {$\sigma_n$ decay};
\end{tikzpicture}%
}
\caption{Finite circular aperture: compact observation operator on far-zone fields;
singular-value decay marks angular harmonics that fall below instrument threshold.}
\label{fig:ch2.signature-aperture-svd}
\end{figure}
Compactness from Subsection~\ref{sec:ch213} yields \(\sigma_n(\mathcal{K})\to 0\); for disk
apertures the decay is approximately exponential in modal index when \(k_0 a_{\mathrm{ap}}\) is
moderate \cite{stratton1941}. In practice the engineer reads off a knee in the singular-value
curve: retain modes above the knee for synthesis or inversion; discard the tail and accept the
residual sidelobe floor as a geometry-limited artifact rather than a tuning defect
\cite{hansen1988}. What practical criterion tells the engineer where to truncate
the spectrum? Retaining every mode is impossible: instrument noise, finite aperture extent, and
finite scan count cap recoverable rank, and singular-value decay controls which source
organizations remain physically realizable through the actual receiver chain
\cite{stratton1941}. Truncation is therefore inevitable; the engineering question is
where the knee lies. Truncating at \(N\) modes with
\begin{equation}
\label{eq:ch2.signature-aperture-truncation}
\frac{\|\mathcal{K}-\mathcal{K}_N\|}{\|\mathcal{K}\|}
\le
\frac{\sigma_{N+1}}{\sigma_1},
\end{equation}
where $\mathcal{K}=\mathcal{O}_{\mathrm{meas}}\circ\mathcal{R}_\omega$ is the full measurement
chain from source to recorded samples, $\mathcal{K}_N$ is its rank-$N$ truncated approximation,
$\sigma_1$ is the largest singular value (dominant export channel), and $\sigma_{N+1}$ is the
first discarded singular value \cite{harrington2001}. The ratio $\sigma_{N+1}/\sigma_1$ is the
relative strength of the strongest discarded radiation channel compared with the dominant
channel---a direct laboratory readout of how much pattern error truncation will introduce
\cite{hansen1988}. Eq.~\eqref{eq:ch2.compression-error-bound} bounds pattern error on \(S_r\). At
\(k_0 a_{\mathrm{ap}}=3\), retaining \(N=9\) singular modes typically captures \(95\%\)
of the radiated power resolvable through the disk; harmonics beyond \(N\approx 12\) lie
below a \(-20\,\mathrm{dB}\) threshold relative to \(\sigma_1\)
\cite{harrington2001}.
Figure~\ref{fig:ch2.signature-aperture-svd} foreshadows Chapter~\ref{ch:04}: realizability
and spectral truncation ask which angular components survive compression by
\(\mathcal{K}\), not which indices a formal expansion can list
\cite{devaney2004}. Laboratory reports therefore inherit a finite rank from
geometry before source inversion is attempted; improving SNR without enlarging
\(a_{\mathrm{ap}}\) cannot recover modes annihilated by \(\sigma_n\) below measurement
noise \cite{hansen1988}.
Local notation ($\sigma_n$, recoverable rank). Compose the measurement operator $\mathcal{K}=\mathcal{O}_{\mathrm{meas}}\circ\mathcal{R}_\omega$
on declared source and field spaces. Its singular values $\sigma_n$---ordered
$\sigma_1\ge\sigma_2\ge\cdots\to 0$---are the radiation efficiencies of independent source
organizations as seen through the actual aperture and receiver chain \cite{stratton1941}. The
\emph{recoverable rank} is the largest $N$ for which $\sigma_N$ remains above the instrument noise
floor; modes beyond that cliff are mathematically listable in a harmonic post-processor yet
physically indistinguishable from clutter on $\Gamma_R$ \cite{chu1948}. Compact
$\sigma_n$ decay and finite dimension on $\mathcal{K}$ are fixed in the display below; elsewhere
invoke Subsection~\ref{sec:ch26.signature-aperture} by reference only \cite{harrington2001}.
The spectral relations developed below restate three laboratory facts in compact form before the eigenbasis is
constructed. A \emph{radiation mode} is a normalized far-zone pattern direction
$\mathbf{u}_n$; its \emph{source mode} $\Jimp^{(n)}$ is an admissible current that excites it.
The \emph{export eigenvalue} $\lambda_n$ of the radiation operator $\mathcal{R}_\omega$---not a
correlation-operator eigenvalue unless later stated---ranks how efficiently that radiation mode
reaches qualified flux on $\Gamma_R$ \cite{stratton1941}. each radiating mode carries export efficiency
$|\lambda_n|$ measured by coupled power in that direction. Second, not every desired far-zone
target can be synthesized: obstructions appear as orthogonality conditions on outgoing
solutions. Third, when geometry possesses rotational symmetry, rotating the hardware reorients
modal coefficients without changing eigenvalue ordering within degenerate blocks. The introductory relation
equations record these facts; the subsections that follow interpret them in measurement
language \cite{stratton1941}.
On a compact range, a \emph{radiation mode} appears as a principal far-zone pattern direction
---a normalized cut on $S_r$ that carries qualified export power---while its \emph{source mode}
is the admissible current distribution that excites it \cite{harrington2001}. Dominant beams
correspond to modes with large $|\lambda_n|$ because those organizations couple most efficiently
to the wattmeter on $\Gamma_R$ \cite{stratton1941}. The eigenpair relation below states the
experimentally familiar observation that certain current organizations radiate much more
effectively than others \cite{collin2001}. On a family of radiating modes, the eigenvalue template reads
\begin{equation}
\label{eq:ch2.modal-eigenvalue-equation}
\mathcal{R}_\omega[\Jimp^{(n)}]=\lambda_n\,\mathbf{u}_n,
\qquad
\mathbf{u}_n=P_n\mathcal{S}_\omega[\Jimp^{(n)}],
\end{equation}
with \(\lambda_n\) measuring radiative export efficiency of mode \(n\) from finite impressed
current \(\Jimp^{(n)}\) \cite{stratton1941}. In synthesis and inversion, $\mathbf{g}$ denotes the
desired far-zone radiated field one wishes to realize; solvability fails when $\mathbf{g}$ carries
a component along cokernel directions---an \emph{orthogonality obstruction} that no admissible
source can remove \cite{stratton1941}. Experimentally, such failure appears as a synthesis target
that refuses to converge, a measured residual orthogonal to every reachable pattern, or a
pattern null that persists despite port tuning \cite{harrington2001}. Solvability on outgoing radiative solutions satisfies
\begin{equation}
\label{eq:ch2.spectral-solvability}
\mathcal{R}_\omega+K
\ \text{solvable on}\ \mathcal{J}_\Omega
\iff
\mathbf{g}\perp\ker(\mathcal{R}_\omega+K)^{\dagger},
\end{equation}
with \(K\) compact on the graph domain of impressed sources \cite{stratton1941}.
When a symmetry group \(G\) commutes with \(\mathcal{R}_\omega\), the equation below expresses
a simple laboratory fact: rotating the hardware and then measuring the pattern gives the same
result as measuring first and rotating the resulting pattern \cite{harrington2001}.
\begin{equation}
\label{eq:ch2.degeneracy-covariance}
\mathcal{R}_\omega\,\mathcal{D}(R)[\Jimp]
=
\mathcal{D}(R)\,\mathcal{R}_\omega[\Jimp],
\qquad R\in G,
\end{equation}
with \(\mathcal{D}(R)\) the field representation induced from
Eq.~\eqref{eq:ch1.active-rotation-action} on \(\mathcal{F}_{\mathrm{rad}}\)
\cite{stratton1941}. Modal projectors compose with radiation selection through
$\Pi_{\mathrm{rad}}$ acting as the radiation filter that passes only exportable content,
$P_n$ as mode-selection lenses that isolate individual beam families from a measured pattern,
and $\Pi_{\mathrm{ev}}$, $\Pi_{\mathrm{non}}$ as projectors onto evanescent and
non-propagating sectors that qualified far-zone meters cannot resolve
\cite{stratton1941}.
\begin{equation}
\label{eq:ch2.spectral-resolution-composition}
\Pi_{\mathrm{rad}}
=
\sum_{n\in\mathcal{N}_{\mathrm{rad}}} P_n,
\qquad
P_n\Pi_{\mathrm{ev}}=P_n\Pi_{\mathrm{non}}=\mathbf{0},
\end{equation}
using Eqs.~\eqref{eq:ch2.spectral-projectors}--\eqref{eq:ch2.rad-null-on-evanescent} by reference
\cite{stratton1941}. Spectral decompositions \emph{organize} export; they do not
\emph{create} angular structure---see Subsection~\ref{sec:ch26.signature-aperture} for the
harmonics-versus-modes distinction stated once in measurement terms \cite{stratton1941}.
